% --- author-revision-loop markup -------------------------------------------
% Paste into the preamble (after \usepackage{xcolor}, or keep the line below).
% \markuptrue  : new text red, deleted text red + struck through (review copy)
% \markupfalse : new text black, deleted text gone (clean copy)
\usepackage{xcolor}
\usepackage[normalem]{ulem}   % normalem keeps \emph as italics
\newif\ifmarkup
\markuptrue
\ifmarkup
  \newcommand{\rn}[1]{{\color{red}#1}}                      % new or changed text
  \newcommand{\rdel}[1]{{\color{red}\sout{#1}}}             % deleted text, same paragraph
  \newcommand{\rdelnote}[1]{{\color{red}\sout{[deleted: #1]}}} % deleted display/figure/table/paragraph
\else
  \newcommand{\rn}[1]{#1}
  \newcommand{\rdel}[1]{}
  \newcommand{\rdelnote}[1]{}
\fi
% Rules for \rdel (ulem):
%  - one paragraph at most; no \par, no display math, no environments inside;
%  - wrap \cite, \ref, \eqref and footnotes in braces: \rdel{as in {\cite{x}}};
%  - inline math is fine: \rdel{the bound $L$ holds}.
%  - replacement: leave a space between old and new, \rdel{an exact} \rn{a}
%    (without it the struck and new words run together: "an exacta").
% Deleting a display, figure, table or several paragraphs: leave a short
% \rdelnote{display (3): definition of K} in place and keep the old source
% commented out between the lines
%   % RDEL-BEGIN <id>
%   % ...old source...
%   % RDEL-END <id>
% so it can be restored. finalize.py removes both.
% ---------------------------------------------------------------------------
